\section{Exercices d'application}

\rubrique{Mode de génération d'une suite}

\exo{}\diff{1}
Soit la suite $(u_n)$ définie pour tout $n\in\N$ par $u_n = \dfrac{2n+1}{n+3}$.
\begin{enumerate}[label=\alph*)]
\item Calculer $u_0$, $u_1$, $u_2$ et $u_{10}$.
\item Exprimer $u_{n+1}$ en fonction de $n$.
\item La suite est-elle définie de façon explicite ou récurrente ? Justifier.
\end{enumerate}

\exo{}\diff{1}
Soit la suite $(v_n)$ définie par $v_0 = 2$ et, pour tout $n\in\N$, $v_{n+1} = 3v_n - 1$.
\begin{enumerate}[label=\alph*)]
\item Calculer $v_1$, $v_2$, $v_3$.
\item Exprimer $v_n$ en fonction de $n$ (on admettra que $v_n = \dfrac{3^n + 1}{2}$).
\item Vérifier cette formule pour $n=0$, $n=1$ et $n=2$.
\end{enumerate}

\rubrique{Sens de variation d'une suite}

\exo{}\diff{1}
Étudier le sens de variation de la suite $(u_n)$ définie par $u_n = n^2 - 4n + 3$ pour $n\in\N$.

\exo{}\diff{2}
Soit $(u_n)$ la suite définie par $u_n = \dfrac{2n+1}{n+2}$ pour $n\in\N$.
\begin{enumerate}[label=\alph*)]
\item Calculer $u_{n+1} - u_n$ et en déduire le sens de variation de $(u_n)$.
\item La suite $(u_n)$ est-elle majorée ? minorée ? bornée ?
\end{enumerate}

\rubrique{Suites arithmétiques}

\exo{}\diff{1}
Une suite arithmétique $(u_n)$ est telle que $u_3 = 8$ et $u_7 = 20$.
\begin{enumerate}[label=\alph*)]
\item Déterminer sa raison $r$ et son premier terme $u_0$.
\item Donner l'expression de $u_n$ en fonction de $n$.
\item Calculer $S = u_0 + u_1 + \cdots + u_{10}$.
\end{enumerate}

\exo{}\diff{2}
Dans une salle de classe du Lycée de Yaoundé, on dispose des chaises en rangées. La première rangée contient 15 chaises, et chaque rangée suivante contient 3 chaises de plus que la précédente. Il y a 12 rangées.
\begin{enumerate}[label=\alph*)]
\item Combien de chaises contient la dernière rangée ?
\item Calculer le nombre total de chaises dans la salle.
\end{enumerate}

\rubrique{Suites géométriques}

\exo{}\diff{1}
Soit $(v_n)$ une suite géométrique de raison $q = 2$ et telle que $v_4 = 48$.
\begin{enumerate}[label=\alph*)]
\item Déterminer $v_0$.
\item Exprimer $v_n$ en fonction de $n$.
\item Calculer $v_1$, $v_2$, $v_3$.
\end{enumerate}

\exo{}\diff{2}
Un capital de 500 000 FCFA est placé à intérêts composés au taux annuel de 4\%. On note $C_n$ le capital disponible au bout de $n$ années.
\begin{enumerate}[label=\alph*)]
\item Exprimer $C_n$ en fonction de $n$.
\item Quel sera le capital au bout de 5 ans ? (arrondir à l'unité)
\item Au bout de combien d'années le capital aura-t-il doublé ?
\end{enumerate}

\rubrique{Suites bornées et applications}

\exo{}\diff{2}
Soit $(u_n)$ la suite définie par $u_0 = 1$ et $u_{n+1} = \dfrac{u_n + 2}{u_n + 1}$ pour $n\in\N$.
\begin{enumerate}[label=\alph*)]
\item Calculer $u_1$, $u_2$, $u_3$.
\item Montrer que pour tout $n\in\N$, $1 \leq u_n \leq 2$.
\item Étudier le sens de variation de $(u_n)$.
\end{enumerate}

\exo{}\diff{3}
Un étudiant de l'Université de Douala économise de l'argent pour un projet. Le premier mois, il dépose 10 000 FCFA, et chaque mois suivant, il dépose 5 000 FCFA de plus que le mois précédent.
\begin{enumerate}[label=\alph*)]
\item Modéliser la situation par une suite $(a_n)$ où $a_n$ est le montant déposé le $n$-ième mois.
\item Calculer le montant déposé au 12\up{e} mois.
\item Quelle est la somme totale économisée au bout d'un an ?
\end{enumerate}

\section{Exercices d'entraînement}

\rubrique{Mode de génération et calculs de termes}

\exo{}\diff{1}
Pour chacune des suites suivantes, calculer $u_0$, $u_1$, $u_2$ et $u_5$ :
\begin{enumerate}[label=\alph*)]
\item $u_n = 3n - 2$
\item $u_n = (-1)^n \dfrac{n}{n+1}$
\item $u_n = \sqrt{n+1} - \sqrt{n}$
\end{enumerate}

\exo{}\diff{1}
Soit $(u_n)$ définie par $u_0 = 1$ et $u_{n+1} = 2u_n + 3$. Calculer $u_1$, $u_2$, $u_3$, $u_4$.

\exo{}\diff{2}
Soit $(u_n)$ définie par $u_0 = 0$, $u_1 = 1$ et $u_{n+2} = u_{n+1} + u_n$ (suite de Fibonacci). Calculer $u_2$, $u_3$, $u_4$, $u_5$.

\rubrique{Sens de variation}

\exo{}\diff{1}
Étudier le sens de variation des suites suivantes :
\begin{enumerate}[label=\alph*)]
\item $u_n = 5 - 2n$
\item $u_n = n^2 + 3n + 1$
\item $u_n = \dfrac{1}{n+1}$
\end{enumerate}

\exo{}\diff{2}
Soit $(u_n)$ définie par $u_n = \dfrac{3n+2}{n+1}$. Montrer que $(u_n)$ est croissante et majorée par 3.

\exo{}\diff{2}
Soit $(u_n)$ définie par $u_0 = 2$ et $u_{n+1} = \dfrac{1}{2}u_n + 1$. Montrer que $(u_n)$ est décroissante et minorée par 2.

\rubrique{Suites arithmétiques}

\exo{}\diff{1}
Une suite arithmétique $(u_n)$ a pour raison $r = -3$ et $u_5 = 7$. Calculer $u_0$ et $u_{10}$.

\exo{}\diff{2}
Soit $(u_n)$ une suite arithmétique telle que $u_2 + u_5 = 20$ et $u_3 + u_7 = 28$. Déterminer $u_0$ et $r$.

\exo{}\diff{2}
Calculer la somme $S = 1 + 4 + 7 + 10 + \cdots + 100$ (termes d'une suite arithmétique).

\exo{}\diff{2}
Dans un stade de football à Limbé, les gradins sont disposés en 20 rangées. La première rangée contient 50 places, et chaque rangée suivante contient 8 places de plus. Calculer le nombre total de places.

\rubrique{Suites géométriques}

\exo{}\diff{1}
Une suite géométrique $(v_n)$ a pour raison $q = \dfrac{1}{2}$ et $v_3 = 4$. Calculer $v_0$ et $v_6$.

\exo{}\diff{2}
Soit $(v_n)$ une suite géométrique telle que $v_2 = 6$ et $v_5 = 48$. Déterminer $q$ et $v_0$.

\exo{}\diff{2}
Calculer la somme $S = 3 + 6 + 12 + 24 + \cdots + 384$.

\exo{}\diff{2}
Un capital de 200 000 FCFA est placé à intérêts composés au taux de 5\% par an. Quelle est la valeur acquise au bout de 8 ans ?

\rubrique{Problèmes mixtes}

\exo{}\diff{3}
Soit $(u_n)$ une suite définie par $u_0 = 1$ et $u_{n+1} = \dfrac{2u_n}{u_n + 1}$.
\begin{enumerate}[label=\alph*)]
\item Calculer $u_1$, $u_2$, $u_3$.
\item Montrer que $(u_n)$ est bornée par 0 et 2.
\item Étudier le sens de variation de $(u_n)$.
\end{enumerate}

\exo{}\diff{3}
Un commerçant de Bafoussam achète un stock de marchandises. Le premier jour, il vend 50 articles, et chaque jour suivant, il vend 10 articles de moins que le jour précédent. Il arrête lorsqu'il ne vend plus rien.
\begin{enumerate}[label=\alph*)]
\item Modéliser par une suite arithmétique.
\item Combien de jours dure la vente ?
\item Combien d'articles a-t-il vendus au total ?
\end{enumerate}

\section{Problèmes type Probatoire/Bac}

\sujetbac[Épargne et placement]
Un jeune entrepreneur de Garoua décide d'épargner pour acheter un terrain. Il place une somme initiale de 100 000 FCFA sur un compte rémunéré à 3\% par an à intérêts composés. De plus, il ajoute 50 000 FCFA à la fin de chaque année (après calcul des intérêts). On note $C_n$ le capital disponible au bout de $n$ années, avec $C_0 = 100\,000$.
\begin{enumerate}[label=\arabic*)]
\item Calculer $C_1$, $C_2$, $C_3$.
\item Montrer que $C_{n+1} = 1,03 C_n + 50\,000$.
\item On pose $D_n = C_n + \dfrac{50\,000}{0,03}$. Montrer que $(D_n)$ est une suite géométrique. En déduire $C_n$ en fonction de $n$.
\item Au bout de combien d'années le capital dépassera-t-il 1 000 000 FCFA ?
\end{enumerate}

\sujetbac[Production agricole]
Une coopérative agricole de l'Ouest-Cameroun produit du maïs. La production de la première année est de 20 tonnes. Chaque année, la production augmente de 2 tonnes par rapport à l'année précédente.
\begin{enumerate}[label=\arabic*)]
\item Modéliser la production annuelle $P_n$ par une suite. Quelle est sa nature ?
\item Calculer la production de la 10\up{e} année.
\item Quelle est la production totale au bout de 10 ans ?
\item La coopérative souhaite atteindre une production annuelle de 50 tonnes. Dans combien d'années cela sera-t-il possible ?
\end{enumerate}

\sujetbac[Suite définie par récurrence]
Soit $(u_n)$ la suite définie par $u_0 = 4$ et $u_{n+1} = \dfrac{1}{2}u_n + 3$ pour $n\in\N$.
\begin{enumerate}[label=\arabic*)]
\item Calculer $u_1$, $u_2$, $u_3$.
\item Montrer que $(u_n)$ est décroissante et minorée par 6.
\item On pose $v_n = u_n - 6$. Montrer que $(v_n)$ est une suite géométrique. En déduire $u_n$ en fonction de $n$.
\item Déterminer la limite de $(u_n)$.
\end{enumerate}

\section{Corrigés}

\corrige{de l'exercice 1}
\begin{enumerate}[label=\alph*)]
\item $u_0 = \dfrac{2\times0+1}{0+3} = \dfrac{1}{3}$ ; $u_1 = \dfrac{3}{4}$ ; $u_2 = \dfrac{5}{5}=1$ ; $u_{10} = \dfrac{21}{13}$.
\item $u_{n+1} = \dfrac{2(n+1)+1}{(n+1)+3} = \dfrac{2n+3}{n+4}$.
\item La suite est définie de façon explicite car $u_n$ est donné directement en fonction de $n$.
\end{enumerate}

\corrige{de l'exercice 2}
\begin{enumerate}[label=\alph*)]
\item $v_1 = 3\times2 - 1 = 5$ ; $v_2 = 3\times5 - 1 = 14$ ; $v_3 = 3\times14 - 1 = 41$.
\item Formule admise : $v_n = \dfrac{3^n + 1}{2}$.
\item Pour $n=0$ : $\dfrac{3^0+1}{2} = \dfrac{2}{2}=1$ (or $v_0=2$, la formule est fausse). Vérifions : $v_0=2$ et $\dfrac{3^0+1}{2}=1$, donc la formule donnée est incorrecte. La bonne formule est $v_n = \dfrac{3^{n+1}+1}{2}$ ? Vérifions : $n=0$ donne $\dfrac{3+1}{2}=2$ ; $n=1$ donne $\dfrac{9+1}{2}=5$ ; $n=2$ donne $\dfrac{27+1}{2}=14$. Donc $v_n = \dfrac{3^{n+1}+1}{2}$.
\end{enumerate}

\corrige{de l'exercice 3}
$u_n = n^2 - 4n + 3$. $u_{n+1} - u_n = (n+1)^2 - 4(n+1) + 3 - (n^2 - 4n + 3) = n^2+2n+1-4n-4+3-n^2+4n-3 = 2n - 3$.
Pour $n=0$, $2\times0-3=-3<0$ ; pour $n=1$, $2-3=-1<0$ ; pour $n=2$, $4-3=1>0$. Donc $(u_n)$ est décroissante pour $n=0,1$ puis croissante à partir de $n=2$.

\corrige{de l'exercice 4}
\begin{enumerate}[label=\alph*)]
\item $u_{n+1} - u_n = \dfrac{2(n+1)+1}{(n+1)+2} - \dfrac{2n+1}{n+2} = \dfrac{2n+3}{n+3} - \dfrac{2n+1}{n+2} = \dfrac{(2n+3)(n+2) - (2n+1)(n+3)}{(n+3)(n+2)}$.
Calculons le numérateur : $(2n+3)(n+2) = 2n^2+4n+3n+6 = 2n^2+7n+6$ ; $(2n+1)(n+3) = 2n^2+6n+n+3 = 2n^2+7n+3$. Différence : $(2n^2+7n+6) - (2n^2+7n+3) = 3$. Donc $u_{n+1} - u_n = \dfrac{3}{(n+3)(n+2)} > 0$. La suite est strictement croissante.
\item $u_n = \dfrac{2n+1}{n+2} = 2 - \dfrac{3}{n+2}$. Comme $\dfrac{3}{n+2} > 0$, $u_n < 2$. De plus $u_n \geq u_0 = \dfrac{1}{2}$. Donc $\dfrac{1}{2} \leq u_n < 2$ : la suite est bornée.
\end{enumerate}

\corrige{de l'exercice 5}
\begin{enumerate}[label=\alph*)]
\item $u_3 = u_0 + 3r = 8$ et $u_7 = u_0 + 7r = 20$. Par soustraction : $4r = 12$ donc $r=3$. Alors $u_0 = 8 - 3\times3 = -1$.
\item $u_n = u_0 + nr = -1 + 3n$.
\item $S = \dfrac{11(u_0 + u_{10})}{2}$ avec $u_{10} = -1 + 30 = 29$. $S = \dfrac{11(-1+29)}{2} = \dfrac{11\times28}{2} = 154$.
\end{enumerate}

\corrige{de l'exercice 6}
Soit $u_n$ le nombre de chaises de la $n$-ième rangée ($n$ commence à 1). $u_1 = 15$, $r=3$.
\begin{enumerate}[label=\alph*)]
\item $u_{12} = u_1 + 11r = 15 + 33 = 48$ chaises.
\item $S = \dfrac{12(u_1 + u_{12})}{2} = \dfrac{12(15+48)}{2} = 6\times63 = 378$ chaises.
\end{enumerate}

\corrige{de l'exercice 7}
\begin{enumerate}[label=\alph*)]
\item $v_4 = v_0 \times 2^4 = 16 v_0 = 48$ donc $v_0 = 3$.
\item $v_n = 3 \times 2^n$.
\item $v_1 = 6$, $v_2 = 12$, $v_3 = 24$.
\end{enumerate}

\corrige{de l'exercice 8}
\begin{enumerate}[label=\alph*)]
\item $C_n = 500\,000 \times (1,04)^n$.
\item $C_5 = 500\,000 \times (1,04)^5 \approx 500\,000 \times 1,21665 = 608\,325$ FCFA.
\item On cherche $n$ tel que $500\,000 \times (1,04)^n \geq 1\,000\,000$, soit $(1,04)^n \geq 2$. $n \geq \dfrac{\ln 2}{\ln 1,04} \approx \dfrac{0,6931}{0,0392} \approx 17,67$. Donc au bout de 18 ans.
\end{enumerate}

\corrige{de l'exercice 9}
\begin{enumerate}[label=\alph*)]
\item $u_1 = \dfrac{1+2}{1+1} = \dfrac{3}{2}=1,5$ ; $u_2 = \dfrac{1,5+2}{1,5+1} = \dfrac{3,5}{2,5}=1,4$ ; $u_3 = \dfrac{1,4+2}{1,4+1} = \dfrac{3,4}{2,4} \approx 1,4167$.
\item Montrons par récurrence que $1 \leq u_n \leq 2$. Initialisation : $u_0=1$ donc vrai. Hérédité : supposons $1 \leq u_n \leq 2$. Alors $u_{n+1} = \dfrac{u_n+2}{u_n+1}$. La fonction $f(x)=\dfrac{x+2}{x+1}$ est décroissante sur $[1,2]$ (dérivée négative). Donc $f(2) \leq u_{n+1} \leq f(1)$, soit $\dfrac{4}{3} \leq u_{n+1} \leq \dfrac{3}{2}$. Donc $1 \leq u_{n+1} \leq 2$. La propriété est héréditaire.
\item $u_{n+1} - u_n = \dfrac{u_n+2}{u_n+1} - u_n = \dfrac{u_n+2 - u_n(u_n+1)}{u_n+1} = \dfrac{u_n+2 - u_n^2 - u_n}{u_n+1} = \dfrac{2 - u_n^2}{u_n+1}$. Comme $u_n \geq 1$, $2 - u_n^2 \leq 1$ ; pour $u_n > \sqrt{2}$, la différence est négative. Calculons : $u_0=1$, $u_1=1,5$, $u_2=1,4$, $u_3\approx1,4167$. La suite semble alterner. Étudions le signe : $u_{n+1} - u_n$ a le signe de $2 - u_n^2$. Donc si $u_n < \sqrt{2}$, la suite croît ; si $u_n > \sqrt{2}$, elle décroît. Elle converge vers $\sqrt{2}$.
\end{enumerate}

\corrige{de l'exercice 10}
\begin{enumerate}[label=\alph*)]
\item $a_n$ est le montant déposé le $n$-ième mois. $a_1 = 10\,000$, $a_{n+1} = a_n + 5\,000$. C'est une suite arithmétique de raison $5\,000$.
\item $a_{12} = a_1 + 11\times5\,000 = 10\,000 + 55\,000 = 65\,000$ FCFA.
\item $S = \dfrac{12(a_1 + a_{12})}{2} = \dfrac{12(10\,000 + 65\,000)}{2} = 6 \times 75\,000 = 450\,000$ FCFA.
\end{enumerate}

\corrige{du problème 1}
\begin{enumerate}[label=\arabic*)]
\item $C_1 = 1,03\times100\,000 + 50\,000 = 103\,000 + 50\,000 = 153\,000$ FCFA. $C_2 = 1,03\times153\,000 + 50\,000 = 157\,590 + 50\,000 = 207\,590$ FCFA. $C_3 = 1,03\times207\,590 + 50\,000 = 213\,817,7 + 50\,000 = 263\,817,7$ FCFA.
\item $C_{n+1} = C_n + 0,03C_n + 50\,000 = 1,03C_n + 50\,000$.
\item $D_n = C_n + \dfrac{50\,000}{0,03} = C_n + \dfrac{5\,000\,000}{3}$. Alors $D_{n+1} = C_{n+1} + \dfrac{5\,000\,000}{3} = 1,03C_n + 50\,000 + \dfrac{5\,000\,000}{3} = 1,03C_n + \dfrac{150\,000 + 5\,000\,000}{3} = 1,03C_n + \dfrac{5\,150\,000}{3} = 1,03\left(C_n + \dfrac{5\,000\,000}{3}\right) = 1,03 D_n$. Donc $(D_n)$ est géométrique de raison $1,03$. $D_0 = 100\,000 + \dfrac{5\,000\,000}{3} = \dfrac{300\,000 + 5\,000\,000}{3} = \dfrac{5\,300\,000}{3}$. Donc $D_n = \dfrac{5\,300\,000}{3} \times (1,03)^n$, et $C_n = D_n - \dfrac{5\,000\,000}{3} = \dfrac{5\,300\,000}{3}(1,03)^n - \dfrac{5\,000\,000}{3}$.
\item On cherche $n$ tel que $C_n \geq 1\,000\,000$. $\dfrac{5\,300\,000}{3}(1,03)^n - \dfrac{5\,000\,000}{3} \geq 1\,000\,000$ $\Rightarrow$ $\dfrac{5\,300\,000}{3}(1,03)^n \geq \dfrac{3\,000\,000 + 5\,000\,000}{3} = \dfrac{8\,000\,000}{3}$ $\Rightarrow$ $(1,03)^n \geq \dfrac{8\,000\,000}{5\,300\,000} = \dfrac{80}{53} \approx 1,5094$. $n \geq \dfrac{\ln 1,5094}{\ln 1,03} \approx \dfrac{0,4118}{0,0296} \approx 13,91$. Donc au bout de 14 ans.
\end{enumerate}

\corrige{du problème 2}
\begin{enumerate}[label=\arabic*)]
\item $P_n$ est la production de la $n$-ième année. $P_1 = 20$, $P_{n+1} = P_n + 2$. C'est une suite arithmétique de raison $2$.
\item $P_{10} = P_1 + 9\times2 = 20 + 18 = 38$ tonnes.
\item $S = \dfrac{10(P_1 + P_{10})}{2} = \dfrac{10(20+38)}{2} = 5\times58 = 290$ tonnes.
\item On cherche $n$ tel que $P_n = 20 + (n-1)\times2 \geq 50$, soit $2(n-1) \geq 30$, $n-1 \geq 15$, $n \geq 16$. Donc dans 15 ans (à partir de la première année), soit la 16\up{e} année.
\end{enumerate}

\corrige{du problème 3}
\begin{enumerate}[label=\arabic*)]
\item $u_1 = \dfrac{1}{2}\times4 + 3 = 2+3=5$ ; $u_2 = \dfrac{1}{2}\times5 + 3 = 2,5+3=5,5$ ; $u_3 = \dfrac{1}{2}\times5,5 + 3 = 2,75+3=5,75$.
\item $u_{n+1} - u_n = \dfrac{1}{2}u_n + 3 - u_n = 3 - \dfrac{1}{2}u_n$. Montrons par récurrence que $u_n \geq 6$. Initialisation : $u_0=4 < 6$, donc la minoration par 6 est fausse. En fait, $u_n$ semble croître vers 6. Vérifions : $u_{n+1} - u_n = \dfrac{6 - u_n}{2}$. Si $u_n < 6$, alors $u_{n+1} > u_n$. Donc $(u_n)$ est croissante. De plus, montrons que $u_n < 6$ pour tout $n$ : par récurrence, $u_0=4<6$ ; si $u_n<6$, alors $u_{n+1} = \dfrac{u_n}{2}+3 < \dfrac{6}{2}+3 = 6$. Donc $(u_n)$ est croissante et majorée par 6.
\item $v_n = u_n - 6$. $v_{n+1} = u_{n+1} - 6 = \dfrac{1}{2}u_n + 3 - 6 = \dfrac{1}{2}u_n - 3 = \dfrac{1}{2}(u_n - 6) = \dfrac{1}{2}v_n$. Donc $(v_n)$ est géométrique de raison $\dfrac{1}{2}$. $v_0 = u_0 - 6 = -2$. Donc $v_n = -2 \times \left(\dfrac{1}{2}\right)^n$, et $u_n = 6 - 2 \times \left(\dfrac{1}{2}\right)^n$.
\item $\lim_{n\to+\infty} \left(\dfrac{1}{2}\right)^n = 0$, donc $\lim u_n = 6$.
\end{enumerate}